\documentclass[11pt]{article}
\usepackage[margin=1in]{geometry}
\usepackage{amsmath,amssymb,amsthm,hyperref}
\newtheorem{lemma}{Lemma}
\newtheorem{theorem}[lemma]{Theorem}
\newtheorem{corollary}[lemma]{Corollary}
\begin{document}
\setlength{\emergencystretch}{1em}
\title{A derivative jet kernel at native prime gaps}
\author{}\date{}\maketitle
We prove a pointwise native-prime offset criterion and the all-order
family $X^{p+1}(A+BX^p)$, with $p$ prime. Its cubic case is support
$\{4,7\}$. A further consequence covers lower exponent four at every
odd prime gap other than eleven, with the stated number-field Gamma
domain. The ordinary quadratic gap is covered separately.
These statements do not assert every positive real shape.

For $d\ge2$, $n\ge1$ and $\gamma>0$ put
\[
 \begin{aligned}
 T_{n,\gamma,d}(t)&=\sum_{k=0}^n\binom nk(\gamma)_{dk}t^k,
       &\epsilon&=(-1)^{d+1},\\
 R&=z^{dn}T_{n,\gamma,d}(\epsilon z^{-d}),
       &M&=\gamma+dn,\\
 E&=M-zD_z, &K_j&=(E)_jR.
 \end{aligned}
\]
Write $B_j=B_{j,\gamma,d}$; the lower shape stays fixed throughout:
\[
 B_0=1,\qquad B_{j+1}=(z+dj+\gamma+d-1)B_j-zB_j',\qquad
 L_j=z^j(D_z-1)^jB_{n+1}.
\]

\begin{lemma}[The upper reduction for any positive offset]
\label{lem:upper}
Let $h\ge1$, $C=M+d$, $A_h=(C)_h$, and
\[
 V_h(X)=\frac{(X)_h-(C)_h}{X-C}
       =\sum_{m=0}^{h-1}(C+m+1)_{h-1-m}(X)_m,\qquad
 \kappa_h=\frac{d(n+1)}{A_h}.
\]
With $K_d=(E)_dR$, the exact identity is
\[
 (\gamma)_h z^{d(n+1)}T_{n+1,\gamma+h,d}(\epsilon z^{-d})
 =\epsilon A_hK_d+
 z^d[A_h-d(n+1)V_h(E)]R.                         \tag{1}
\]
At $R=0$, the upper zero is therefore equivalent to
$K_d=\epsilon\kappa_h z^dV_h(E)R$.
\end{lemma}
\begin{proof}
The finite telescoping identity for the displayed coefficients follows
by subtracting successive terms
$(X)_m(C+m)_{h-m}$.
Pascal and Gamma shifting give the upper operator
$z^d(E)_h+\epsilon(E)_{d+h}$.
The lower Euler equation gives
\[
 (E-M)K_d=-\epsilon z^d(E-\gamma)R,\qquad
 Ez^d=z^d(E-d).
\]
For $F(E)=(E+d)_h$, polynomial division gives
$F(E)-F(M)=(E-M)G(E)$ with $G(E-d)=V_h(E)$.
Consequently
\[
 F(E)K_d=A_hK_d-\epsilon z^dV_h(E)(E-\gamma)R.
\]
Use $(E)_{d+h}=F(E)(E)_d$ and
$(E)_h=A_h+(E-C)V_h(E)$ to obtain (1).
Every step is a finite operator identity, with no condition $h\le d$.
\end{proof}

The necessary concomitant input is the proof of
\texttt{thm:mixed}, equation (3), in the
\href{LOWER_TWO_PRIME_GAP_SFC.tex}{lower-two prime-gap note}.
After its nonzero gauges are removed, at $R=0$ it says
\[
 B_n(\xi z)K_d-
 \sum_{j=0}^{d-2}L_j(\xi z)K_{d-1-j}=0,\qquad \xi^d=1. \tag{2}
\]
The full concomitant has rank $d$ on the $(d+1)$-jet and kernel
the actual nonzero $R$ jet. Fourier evaluation is used only in
characteristic zero, giving a cyclic coefficient matrix $Z(u)$
on $(K_1,\ldots,K_d)$ at $u=z^d$:
its last column is $\mathcal C_d(B_n;u)$ and its column $m<d$
is $-\mathcal C_d(L_{d-1-m};u)$.
The identities transfer algebraically to characteristic-zero fields.
No evaluation of an exponential gauge in a $p$-adic field is needed.
The upper-operator calculation also appears in the proof of
\texttt{thm:offset-rank} of the
\href{DIAGONAL_PRIME_GAP_SFC.tex}{diagonal-prime note};
the lemma above supplies the required offset beyond that theorem's range.

\begin{lemma}[The reduced rising-jet kernel]\label{lem:kernel}
Let $p$ be prime, let $\gamma$ be integral in a complete algebraically
closed extension of $\mathbb Q_p$, and let $u$ be an integral root of
the lower monic reciprocal. The reduction of $Z(u)$ at $d=p$
has rank $p-1$ and kernel generated by
\[
 k_m=\left.\frac{d}{dx}(x)_m\right|_{x=\bar\gamma},
 \qquad1\le m\le p,\qquad k_p=-1.
                                                               \tag{3}
\]
No condition $\bar\gamma\in\mathbb F_p$ is imposed.
\end{lemma}
\begin{proof}
The monic reciprocal is integral, and
\[
 (\gamma)_{pk}\bmod\mathfrak m
     =(\bar\gamma^p-\bar\gamma)^k,\qquad
 u^nT_{n,\gamma,p}(\epsilon/u)\bmod\mathfrak m
     =(u+\bar\gamma^p-\bar\gamma)^n.
\]
Here $\bar\epsilon=1$ at the native prime.
Thus $\bar u+\bar\gamma^p-\bar\gamma=0$.
On the reduced cyclic space $k[z]/(z^p-\bar u)$,
\[
 A=z-zD_z+\bar\gamma-1,\qquad A^p=A.
\]
Indeed its cyclic bidiagonal characteristic polynomial is
\[
 X^p-X-(\bar u+\bar\gamma^p-\bar\gamma)=X^p-X.
\]
Its first $p$ powers of $1$ are monic in successive degrees, so $1$
is cyclic. The distinct eigenvalues are $\mathbb F_p$.
The reduced auxiliary and concomitant columns are
\[
 B_n=A^n1,\qquad
 L_j=(-1)^j A^{n+1}(A-\bar\gamma+1)_j1.
\]
The second identity follows from
$L_{j+1}=z(D_z-1)L_j-jL_j$.
These $L_j$, $0\le j\le p-2$, are a unit triangular combination of
$A^{n+1}1,\ldots,A^{n+p-1}1$.
The nonzero eigenspaces show that their rank is $p-1$.
All columns have zero component in the zero eigenspace because $n\ge1$.

It remains to check (3). At a nonzero eigenvalue $\lambda\in\mathbb F_p$,
the required equality is
\[
 k_p=\sum_{j=0}^{p-2}(-1)^j\lambda
          (\lambda-\bar\gamma+1)_j k_{p-1-j}.
\]
Differentiate the rising Vandermonde identity
\[
 (x+y)_{p-1}=\sum_{j=0}^{p-1}\binom{p-1}{j}
                     (y)_j(x)_{p-1-j}
\]
in $x$ with $y$ fixed, then set
$x=\bar\gamma$, $y=\lambda-\bar\gamma+1$.
Since $\binom{p-1}{j}=(-1)^j$ in characteristic $p$, the sum
is $\lambda\,((x)_{p-1})'|_{x=\lambda+1}$.
Now $(x)_{p-1}=(x^p-x)/(x-1)$ and
$\lambda+1-1=\lambda\ne0$, so this value is $-1=k_p$.
The omitted term at $j=p-1$ has derivative of $(x)_0$, hence is zero.
Rank $p-1$ makes this the whole kernel.
\end{proof}

\begin{theorem}[A pointwise native-prime offset criterion]
\label{thm:criterion}
Let $p$ be prime, $1\le h\le p+1$, $n\ge1$, and let $\gamma>0$
be real algebraic with an embedding
$\mathbb Q(\gamma)\longrightarrow\mathbb C_p$ having integral image.
In this embedding put
\[
 C=\gamma+pn+p,\quad \kappa_h=\frac{p(n+1)}{(C)_h},\qquad
 \sigma_h(x)=[Y^2](x+Y)_h .
\]
If either
\[
 \sigma_h(\bar\gamma)\ne0
 \qquad\hbox{or}\qquad v_p(\kappa_h)\ge0,                 \tag{4}
\]
then $T_{n,\gamma,p}$ and $T_{n+1,\gamma+h,p}$ are coprime over
$\mathbb Q(\gamma)$ and over $\mathbb C$.
The notation $\sigma_h$ is the second Hasse derivative; it does
not divide an ordinary derivative by two.
\end{theorem}
\begin{proof}
Use the chosen embedding and let a common root give
$u=\epsilon/t$. Its monic reciprocal has integral coefficients,
so $u$ is integral. It is nonzero.
The actual nonzero rising jet at $R=0$ can be divided by an actual
component of minimum valuation to give a primitive integral vector
$w=(w_1,\ldots,w_p)$.
The full jet cannot vanish at a nonzero ordinary point of the
monic order-$p+1$ lower equation.
Lemma~\ref{lem:kernel} gives $\bar w=a k$, $a\ne0$, so $w_p$ is a unit.

Let $v=(v_1,\ldots,v_p)$ be the coefficients of $V_h$ in the rising
basis, with zeros above its degree. All are integral because
$h\le p+1$ and $C$ is integral.
The constant term multiplies $R=0$ and is omitted.
Evaluation on the kernel gives
\[
 \sum_{m=1}^pv_m k_m=V_h'(\bar\gamma)=\sigma_h(\bar\gamma).
                                                               \tag{5}
\]
Indeed $\bar C=\bar\gamma$, and expansion at $X=C$ of
$((X)_h-(C)_h)/(X-C)$ has linear coefficient $\sigma_h(C)$.

Set $\eta=\epsilon\kappa_hu$. The upper condition (1) is
$w_p=\eta\,v w$.
If $v_p(\kappa_h)\ge0$ and $\bar\gamma\in\mathbb F_p$,
the reciprocal reduction gives $u\in\mathfrak m$ and thus
$v(\eta)>0$, a contradiction.
If $\bar\gamma\notin\mathbb F_p$, every factor of $(C)_h$
is a unit, because $h\le p+1$ and its factors reduce to
$\bar\gamma,\bar\gamma+1,\ldots$.
Then $v(\kappa_h)>0$ and the same contradiction follows
with merely integral $u$.

It remains that $\sigma_h(\bar\gamma)\ne0$ and $v(\kappa_h)<0$.
Equation (5) makes $v w$ a unit; hence $\eta$ is a unit.
Now $q=(\epsilon\kappa_h)^{-1}$ lies in $\mathfrak m$,
and $u=q\eta$ has positive valuation.
The reciprocal reduction now gives
$\bar\gamma^p-\bar\gamma=0$, hence $\bar\gamma\in\mathbb F_p$.
Choose $p-1$ independent rows $I$ of the reduction $Z(0)$ and
form the exact finite determinant
\[
 F(X)=\det\begin{bmatrix}
          Z_I(qX)\\ e_p^{\mathsf T}-Xv
        \end{bmatrix}.                                  \tag{6}
\]
It has integral coefficients.
Its reduction is affine: the first $p-1$ rows are constant since
$q\in\mathfrak m$.
The row kernel is spanned by $k$, with $k_p=-1$ and $v k=\sigma_h$.
Thus $\bar F(0)\ne0$ and the coefficient of $X$ is nonzero;
equivalently $\bar F=c(1+\sigma_h X)$ for a unit $c$.
Any finite integral polynomial with affine nonzero-slope reduction
is injective on the entire valuation ring: its divided difference
is a unit. This remains true when the valuation of $q$ is fractional.

Every actual common root gives an integral zero of this same $F$.
The row set $I$ was chosen at $u=0$, independent of the root.
The map $t\mapsto\eta=\kappa_h/t$ is injective.
There is at most one distinct common root over $\mathbb C_p$.
Every conjugate over the fixed coefficient field $\mathbb Q(\gamma)$
of an actual common factor obeys the same identities.
Separability forces that factor to have degree one.
Its root is real in the original real embedding, contradicting
strict positivity of the even-index member:
\[
 T_{j,\alpha,p}(t)=\frac1{\Gamma(\alpha)}
   \int_0^\infty x^{\alpha-1}(1+tx^p)^j e^{-x}\,dx>0
 \quad(j\text{ even},\ \alpha>0,\ t\in\mathbb R).
\]
Here the relevant shape is $\gamma$ or $\gamma+h$.
A Bezout identity over the coefficient field gives the complex claim.
This reasoning uses all conjugates, not rationality of an arbitrary
unique integral zero of the auxiliary determinant.
\end{proof}

\begin{theorem}[Prime plus one lower exponent]\label{thm:family}
Let $p$ be prime and let $\beta>0$ be real algebraic.
Suppose some embedding $\mathbb Q(\beta)\longrightarrow\mathbb C_p$
has $v_p(\beta)\ge0$.
Every nonzero complex phase
\[
                         X^{p+1}(A+BX^p)
\]
has a nonzero Gamma-$\beta$ moment in every consecutive positive pair.
\end{theorem}
\begin{proof}
For odd $p$, in characteristic $p$,
\[
 (x)_{p+1}=x^{p+1}-x^2,\qquad
 \sigma_{p+1}(x)=\binom{p+1}{2}x^{p-1}-1=-1.
\]
Apply Theorem~\ref{thm:criterion} with
$\gamma=(p+1)n+\beta$ and $h=p+1$.
The exact phase normalization is
\[
 \mathcal L_\beta((X^{p+1}(A+BX^p))^j)
 =A^j(\beta)_{(p+1)j}T_{j,(p+1)j+\beta,p}(B/A).
\]
The nonzero factors do not affect zeros.
For $p=2$, the phase is $X^3(A+BX^2)$, covered at every positive
real shape by
\texttt{cor:gap-two-}\allowbreak\texttt{lower-two-three} of the
\href{lower_exponent_initial_and_boundary.tex}{lower-exponent note}. If either coefficient vanishes, the nonzero
phase is a monomial and all positive Gamma moments are nonzero.
\end{proof}

\begin{corollary}[Lower exponent four at almost every prime gap]
\label{cor:lower-four-prime}
Let $p$ be an odd prime other than eleven, and let $\beta>0$ be
real algebraic with an integral embedding at one $p$-adic place.
Every nonzero $X^4(A+BX^p)$ has a nonzero Gamma-$\beta$ moment in
every positive consecutive pair.
At ordinary shape one, the same conclusion also holds at prime
two, by the quadratic-gap theorem.
No all-order assertion at gap eleven is made here.
\end{corollary}
\begin{proof}
For $h=4$,
\[
 \sigma_4(x)=6x^2+18x+11,\qquad
 (\sigma_4(0),\sigma_4(-1),\sigma_4(-2),\sigma_4(-3))
       =(11,-1,-1,11).
\]
Use $\gamma=4n+\beta$ in Theorem~\ref{thm:criterion}.
If $(C)_4$ is a unit, then $\kappa_4\in\mathfrak m$.
Otherwise $\bar\gamma$ is one of $0,-1,-2,-3$, where the displayed
values of $\sigma_4$ are units for the stated primes.
Thus the criterion always applies.
For the ordinary prime-two case use
\texttt{cor:ordinary-support-}\allowbreak\texttt{four-six} of the
\href{LOWER_FOUR_QUADRATIC_GAP_SFC.tex}{lower-four quadratic-gap note}.
\end{proof}

\begin{corollary}[The cubic-gap support four/seven]\label{cor:four-seven}
Every ordinary nonzero $X^4(A+BX^3)$ satisfies Strong Factorial at
every positive consecutive pair. The same holds at each positive
real algebraic Gamma shape integral at one three-adic place.
\end{corollary}
\begin{proof}
Use Theorem~\ref{thm:family} at $p=3$, including $\beta=1$.
\end{proof}

\begin{lemma}[The explicit cubic offset-four rank form]
\label{lem:cubic-offset-four}
For arbitrary $\gamma>0$, $n\ge1$ and $u\ne0$, put
\[
 M=\gamma+3n,\quad \kappa=\frac{3(n+1)}{(M+3)_4},\quad
 \eta=\kappa u,\quad B=B_{n,\gamma,3},\
 C=B_{n+1,\gamma,3},\ D=B_{n+2,\gamma,3}.
\]
The two equalities
$T_{n,\gamma,3}(1/u)=T_{n+1,\gamma+4,3}(1/u)=0$
hold if and only if the cyclic columns of
\[
                 [(1-\eta)C-\eta(M+6)B,\ D]              \tag{7}
\]
have rank less than two. This includes $\eta=1$.
\end{lemma}
\begin{proof}
At $R=0$, set $J=K_2+(M+5)K_1$.
Lemma~\ref{lem:upper} and the concomitant give
\[
 B K_3-CJ+DK_1=0,\qquad
 (1-\eta)K_3=\eta(M+6)J.
\]
Since $\eta(M+6)\ne0$, eliminate $J$ to obtain (7);
no division by $1-\eta$ is used.
At a common root, $K_1=K_3=0$ would force $J=K_2=0$,
contradicting the nonzero full jet.
Conversely a nonzero nullvector in $(K_1,K_3)$ determines $J$,
then $K_2$. The full concomitant kernel identifies this jet with
the actual $R$ jet and supplies both restrictions.
Fourier evaluation converts these identities to cyclic coefficient
columns in characteristic zero. This proves both directions.
\end{proof}

The criteria above do not cover every Gamma shape or arbitrary supports.
The written identities and valuation argument supply their unbounded
quantifiers; a finite replay cannot replace those arguments.
\end{document}
